% !TeX root = ./main.tex

\chapter{绪论}\label{ch:intro}

\section{研究背景与意义}\label{sec:background}

随着大语言模型（Large Language Models, LLMs）的飞速发展，其在自然语言理解与生成任务上展现出卓越的能力。然而，纯参数化的 LLM 面临知识时效性不足、事实性偏差（幻觉）以及难以追溯推理依据等问题。检索增强生成（Retrieval-Augmented Generation, RAG）技术通过在生成回答前从外部知识库检索相关证据，将参数化知识与非参数化外部知识相结合，有效缓解了上述问题~\cite{lewis2020rag}。RAG 的核心价值在于提升知识密集型任务的事实性、可更新性和证据可追踪性，已在企业知识问答、学术文献分析、法律文书检索等场景得到广泛应用。

然而，真实世界中的文档远非纯文本集合。学术论文、企业财报、技术手册、合同文书和教学课件等文档通常同时包含正文、公式、表格、图表、图像、图注、脚注、跨页段落和图表引用等多种信息形态。传统文本 RAG 主要依赖 OCR 或 PDF 文本抽取后的纯文本向量相似度匹配，在处理这类图文混排的复杂文档时面临三个根本性挑战：

\begin{enumerate}
\item \textbf{视觉信息丢失}：图表趋势、版面结构、颜色标记、坐标轴和表格边界等视觉信息很难被纯文本准确保留。许多关键语义（如折线图的趋势方向、柱状图的对比关系、流程图的步骤先后）天然依赖视觉表达，仅靠文本描述难以完整捕捉。
\item \textbf{结构关系不足}：同一页内的图文对应关系、跨页段落的延续关系、图号与结论之间的支撑关系、关键词与内容片段的锚定关系等，很难仅通过向量相似度稳定表达。真实文档中的重要信息往往不是线性排列，而是通过复杂的引用关系相互关联。
\item \textbf{证据链不可解释}：向量检索擅长找到与查询语义相似的孤立文本片段，但不擅长回答``某结论由哪些图表和实验支撑''这类多跳推理问题。用户难以追溯答案的证据来源和推理路径。
\end{enumerate}

GraphRAG 技术通过在传统 RAG 的基础上引入图结构，为解决上述问题提供了新的思路。GraphRAG 将知识库中的文本片段、实体、概念、关系和事件组织为知识图谱，使检索不仅依赖语义相似度，还能沿图谱中的关系边进行多跳扩展~\cite{edge2024graphrag}。对于文档问答而言，图结构可以将页面、图表、表格、段落、关键词和结论等异质元素统一组织到同一证据空间中，通过显式的关系边捕获跨页延续、图文支撑和多跳推理路径。

与此同时，视觉语言模型（Vision-Language Models, VLMs）的快速发展为多模态文档理解提供了强大的技术基础。以 Qwen-VL 系列为代表的视觉语言模型能够直接接收页面图像和文本提示，完成 OCR 识别、版面理解、图表描述、结构化抽取和视觉问答等任务~\cite{qwen2025qwen3vl}。这为构建端到端的多模态文档问答系统提供了便利：不再需要组合 OCR 引擎、版面分析器、表格识别器等多个独立模块，而是由一个统一的 VLM 完成多模态文档理解的各项子任务。

\section{研究目标与内容}\label{sec:objectives}

本项目面向中山大学 2025--2026 春季学期``多模态大模型原理与应用''课程题目 3``文档理解 + 多模态检索问答系统''，目标是设计并实现一个能够处理 PDF、扫描文档、论文页面和演示文稿截图的多模态 GraphRAG 文档问答系统。

本研究的核心问题是：

\begin{quote}
对于图文混排、跨页引用和结构复杂的长文档，如何把页面视觉信息、文本语义片段和知识图谱关系融合到同一条检索问答链路中，使回答既能利用多模态证据，又能给出可追溯的页码、节点和关系依据？
\end{quote}

围绕这一核心问题，本项目的具体研究内容包括：

\begin{enumerate}
\item \textbf{多模态文档解析与结构化}：设计基于 VLM 的页面图像结构化抽取流水线，从文档页面中同时提取文本块、表格、图表描述、结论、关键词和语义关系，生成结构化的多模态文档表示。
\item \textbf{异构图谱构建}：设计面向文档问答的异构文档图结构，将文本、表格、图表、结论、关键词等内容节点通过同页边、关键词边、跨页延续边和语义关系边组织为统一的知识图谱。
\item \textbf{GraphRAG 检索与生成}：设计融合向量检索、图谱邻域扩展和社区级全局补充的多阶段检索策略，以及将文本证据、关系路径和页面图像融合为多模态上下文的答案生成方法。
\item \textbf{系统实现与前端可视化}：实现完整的端到端系统原型，包括文档入库、知识库管理、在线问答、图谱可视化等核心功能，提供可交互的 Web 界面。
\item \textbf{实验评测与消融分析}：构建自建测试集，设计多组消融实验，系统性评估向量检索、知识图谱和多模态证据各组件对问答性能的贡献。
\end{enumerate}

\section{论文组织结构}\label{sec:organization}

本文其余章节的组织如下：

\begin{itemize}
\item \textbf{第二章 相关工作}：系统回顾 RAG、GraphRAG、多模态文档理解和视觉语言模型等领域的研究现状，分析现有工作的优势与不足，明确本项目的技术定位。
\item \textbf{第三章 系统设计与实现}：详细阐述系统的总体架构、核心模块设计、关键技术实现细节，包括文档解析流水线、图谱构建算法、GraphRAG 检索与生成策略、后端 API 和前端界面。
\item \textbf{第四章 实验设计与结果分析}：介绍自建测试集的构建过程、评测指标设计、消融实验方案，报告并分析五组消融实验的定量结果，讨论不同组件对系统性能的影响。
\item \textbf{第五章 总结与展望}：总结本文的主要工作和贡献，分析系统的局限性，展望未来改进方向。
\end{itemize}
